\section*{ЗАКЛЮЧЕНИЕ}
\phantomsection
\addcontentsline{toc}{section}{Заключение}

В ходе выполнения курсовой работы проведён численный анализ влияния дискретного микрорельефа на трибологические характеристики радиального подшипника скольжения. Решены следующие задачи:

\begin{enumerate}
    \item Выполнен обзор теоретических основ гидродинамической теории смазки и современных работ, посвящённых лазерному текстурированию поверхностей трения. Показано, что основным механизмом улучшения характеристик является микрогидродинамический эффект, возникающий на локальных клиновых зазорах при обтекании кромок углублений.

    \item Изучена методика численного решения уравнения Рейнольдса для подшипника конечной длины с использованием метода конечных разностей, итерационной процедуры Гаусса-Зейделя и верхней релаксации. Расчёт выполнен на сетке $500 \times 500$ узлов при критерии сходимости $\delta < 10^{-5}$.

    \item Рассчитаны поля давления и интегральные характеристики (нагрузочная способность $F$, коэффициент трения $\mu$, расход смазки $Q$) для гладкого подшипника и подшипника с эллипсоидальными углублениями (тип~1) при эксцентриситетах $\varepsilon = 0{,}05 - 0{,}80$. Геометрия подшипника: $R = 50$~мм, $c = 70$~мкм, $L = 80$~мм.

    \item Проведён сравнительный анализ гладкого и текстурированного подшипников. При рабочем эксцентриситете $\varepsilon = 0{,}6$ получены следующие результаты:
    \begin{itemize}
        \item нагрузочная способность возросла на $4{,}59$\%: с $F = 10685{,}6$~Н до $F = 11175{,}6$~Н;
        \item коэффициент трения уменьшился на $5{,}21$\%: с $\mu = 0{,}0151$ до $\mu = 0{,}0143$;
        \item расход смазки увеличился на $0{,}51$\%: с $Q = 0{,}6739$~л/с до $Q = 0{,}6774$~л/с.
    \end{itemize}
\end{enumerate}

Полученные данные свидетельствуют о том, что эллипсоидальные углубления с параметрами $a = 2{,}41$~мм, $b = 2{,}214$~мм, $h_p = 10$~мкм ($h_p/c = 0{,}1$) оказывают положительное влияние на рабочие характеристики подшипника геометрии~B. Повышение нагрузочной способности при одновременном снижении коэффициента трения и незначительном росте расхода подтверждает эффективность данного типа микрорельефа для подшипников увеличенного типоразмера.

Плавный профиль эллипсоидальных углублений обеспечивает постепенное изменение зазора от центра к периферии углубления, исключая резкие скачки давления и разрывы потока. Это делает эллипсоидальный микрорельеф перспективным для применения в подшипниках, эксплуатируемых при стационарных нагрузках и умеренно-высоких эксцентриситетах.

Проведённый анализ позволяет рекомендовать использование эллипсоидального текстурирования для подшипников скольжения геометрии~B в нефтегазовом оборудовании, функционирующем в условиях жидкостного трения. Наибольший эффект достигается в тяжелонагруженных режимах ($\varepsilon > 0{,}5$), где прирост несущей способности и снижение трения проявляются в наибольшей степени.

